\documentclass[a4paper,12pt,french]{report}

% =======================
% IMPORTATION DES PACKAGES
% =======================

\usepackage[T1]{fontenc}
\usepackage[utf8]{inputenc}
\usepackage[margin=2cm]{geometry}
\usepackage{theorem}
\usepackage{lmodern}
\usepackage{amssymb, amsmath, amsfonts, dsfont}
\usepackage{babel}

\renewcommand{\familydefault}{\sfdefault}

\pagestyle{headings}

\theoremstyle{break}
\theoremheaderfont{\scshape}
\theorembodyfont{\upshape}

\newtheorem{defin}{Définition}[section]
\newtheorem{theo}[defin]{Théorème}
\newtheorem{prop}[defin]{Proposition}

\title{MÉSURE ET INTÉGRATION}
\date{\today}
\author{}

\begin{document}
\maketitle
\textsc{Précision} \\
Ceci est juste des notes pour m'aider à comprendre la théorie de la mesure et intégration tout en apprenant latex.
Mais si cela peut vous aider, vous pouvez le consulter et le télécharger.

\tableofcontents

\chapter{THÉORIE DES ENSEMBLES}

\section{Théorie des ensembles}
\begin{defin}
Soit $X$ un ensemble non vide et $(A_{i})_{i \in \mathbb{N}}.$ 
      $$\bigcup_{i \in \mathbb{N}} A_{i} := \{ x \in X | \mbox{ il existe } i_{0} \mbox{ tel que } x \in A_{i_{0}}\}.$$
      $$\bigcap_{i \in \mathbb{N}} A_{i} := \{ x \in X | \mbox{ pour tout } i_{0} \mbox{ tel que } x \in A_{i_{0}}\}.$$
      $$A \backslash B := \{ x \in A \mbox{ et } x \not\in B\}.$$
      $$X\backslash A := \{ x \not\in A\}.$$
      $$ A \triangle B := (A \backslash  B) \cap (B \backslash A).$$ 
\end{defin}

\begin{prop}[Loi de Morgan]
     $$X\backslash \left(\bigcup_{i \in \mathbb{N}} A_{i}\right) = \bigcap_{i \in \mathbb{N}} (X \backslash A_{i}).$$
     $$X\backslash \left(\bigcap_{i \in \mathbb{N}} A_{i}\right) = \bigcup_{i \in \mathbb{N}} (X \backslash A_{i}).$$
\end{prop}

\begin{prop}[Distributivité]
     $$A \cap (B \cup C) = (A \cap B) \cup (A \cap C).$$
     $$A \cup (B \cap C) = (A \cup B) \cap (A \cup C).$$
\end{prop}

\begin{defin}
       $$\mathds{1}_{A}(x) = \left\{ \begin{array}{c} 1 \mbox{ si } x \in A \\ 0 \mbox{ si } x \not\in A\end{array}\right.$$
\end{defin}

\begin{prop}
       \begin{enumerate}
              \item $\mathds{1}_{X\backslash A} = 1 - \mathds{1}_{A}$,
              \item $\mathds{1}_{A \cap B} = \mathds{1}_{A} \mathds{1}_{B}$,
              \item $\mathds{1}_{A \cup B} = \mathds{1}_{A} + \mathds{1}_{B} - \mathds{1}_{A \cap B}$,
              \item $\mathds{1}_{A \backslash B} = \mathds{1}_{A} - \mathds{1}_{A \cap B}$,
              \item $\mathds{1}_{A \triangle B} = \mathds{1}_{A} + \mathds{1}_{B} \mbox{ mod } 2$.
       \end{enumerate}
\end{prop}

\section{Ensemble ordonnée}
\begin{defin}[Relation d'ordre]
       Soit X un ensemble non vide. On dit que  $\preccurlyeq$ est une relation d'ordre s'il verifie :
       \begin{itemize}
              \item[i)] Reflexivité : Soit $x \in X$, $x \preccurlyeq x$,
              \item[ii)] Anti-symetrie : Soit $x, y \in X$, $x \preccurlyeq y$ et $y \preccurlyeq x$, alors $x=y$,
              \item[iii)] Transitivité : Soit $x, y, z \in X$, $x \preccurlyeq y$ et $y \preccurlyeq z$, on a $x \preccurlyeq z$.   
       \end{itemize}
       On appelle $(X, \preccurlyeq)$ un ensemble ordonnée.
\end{defin}

\begin{defin}[Relation d'ordre total et partielle]
       On dit que l'ensemble ordonnée $(X, \preccurlyeq)$ est total si tout élement de X est comparable c'est à dire pour $x, y \in X$, $x \preccurlyeq y$ ou bien $y \preccurlyeq x$.
       Si c'est le contraire, on dit qu'elle est partielle.
\end{defin}

\chapter{TRIBUS, CLAN ET CLASSE MONOTONE}
\section{Terminologie}
\begin{defin}[Clan]
      Soit X un ensemble non vide. On dit que $\mathcal{A}$ est un clan sur X si elle verifie :
      \begin{itemize}
             \item $\emptyset \in \mathcal{A}$,
             \item $\mathcal{A}$ est stable par passage au complémentaire, c'est à dire pour $A \in \mathcal{A}$, on a $X \backslash A \in \mathcal{A}$,
             \item $\mathcal{A}$ est stable par réunion, c'est à dire pour $A, B \in \mathcal{A}$, on a $A \cup B \in \mathcal{A}$.
      \end{itemize}
\end{defin}

\begin{defin}[Tribus]
      Soit X un ensemble non vide. On dit que $\mathcal{A}$ est une tribu sur X si elle vérifie :
      \begin{itemize}
             \item $\emptyset \in \mathcal{A}$,
             \item $\mathcal{A}$ est stable par passage au complémentaire, c'est à dire pour $A \in \mathcal{A}$, on a $X \backslash A \in \mathcal{A}$,
             \item $\mathcal{A}$ est stable par réunion dénombrable, c'est à dire pour $(A_{n})_{n \in \mathbb{N}} \in \mathcal{A}$, on a $\displaystyle \bigcup_{n \in \mathcal{N}} A_{n} \in \mathcal{A}$.
      \end{itemize}
\end{defin}

\begin{defin}[Classe monotone]
       On dit que $\mathcal{A}$ est une classe monotone si elle vérifie : 
       \begin{itemize}
              \item $(A_{n})_{n \in \mathbb{N}} \nearrow A$, alors $\displaystyle\bigcup_{n \in \mathbb{N}} A_{n} \in \mathcal{A}$
              \item $(A_{n})_{n \in \mathbb{N}} \searrow A$, alors $\displaystyle\bigcap_{n \in \mathbb{N}} A_{n} \in \mathcal{A}$
       \end{itemize}
\end{defin}

\chapter{MESURABILITÉ}
\section{Espace mesurable}
\begin{defin}
       Soit X un ensemble non vide. On dit que $(X,\mathcal{A})$ est un espace mesurable si $\mathcal{A}$ est un tribus.
\end{defin}
\end{document}